% Source: https://artofproblemsolving.com/community/c3h1753584_geometry_help_please
% Stable identity: aops-topic-1753584-q01
Let $M$ be the midpoint of $AB$.

\begin{asy}
usepackage("amsmath");
size(5cm);
defaultpen(fontsize(10pt));
pair O=(0,0),A=6*dir(30),B=6*dir(150);
fill(arc(O,6,30,150)--cycle,paleblue);
draw(circle(O,6),linewidth(0.8));
draw(O--A,linewidth(0.8));draw(O--B,linewidth(0.8));draw(A--B,linewidth(0.8));
draw(arc(O,1,30,150),linewidth(0.6));label("$120^\circ$",O+1.75*dir(118),fontsize(8pt));
pair M=(A+B)/2;draw(O--M,dashed);
draw(M--M+(0.4,0)--M+(0.4,-0.4)--M+(0,-0.4));
dot(O);label("$O$",O,S);label("$A$",A,E);label("$B$",B,W);label("$M$",M,dir(60),fontsize(8pt));
label("$6$",(O+A)/2,SE);
label("\mbox{}",(min(currentpicture,user=true)+max(currentpicture,user=true))/2,black); // keeps the text after the figure black when LaTeX typesets the labels inline
\end{asy}

The sector has area $\dfrac{120}{360}\cdot\pi\cdot6^2=12\pi$. In the isosceles triangle $AOB$, the altitude $OM$ bisects the $120^\circ$ angle, so triangle $OMA$ is a $30^\circ$--$60^\circ$--$90^\circ$ triangle with $OM=3$ and $AM=3\sqrt3$. So $AB=6\sqrt3$ and the triangle has area $\dfrac{1}{2}\cdot6\sqrt3\cdot3=9\sqrt3$. (Equivalently, $\dfrac{1}{2}\cdot6\cdot6\sin120^\circ=9\sqrt3$.) The shaded region has area
\[
\boxed{12\pi-9\sqrt3}\approx22.11\text{ cm}^2 .
\]
